\section{Finite-field models and an explicit eventual range}
\label{sec:finite-field}

This section constructs an ordinary-index-zero ample bundle on every mixed surface with a finite-field model, and bounds the ample difference needed for the sharp rank formula.

We use two elementary consequences of Frobenius and Verschiebung. If $G$ is a finite connected group scheme of order $h$ over $\FF_q$, its coordinate algebra is local with augmentation ideal $I$ and $I^h=0$. Indeed, every nonzero term in the sequence $I,I^2,\ldots$ strictly decreases dimension until it is zero: equality of two successive terms would imply their vanishing by Nakayama's lemma. Thus the $q^m$-power Frobenius kills $G$ whenever $q^m\ge h$. If the Cartier dual $G^D$ is also connected, the same argument on $G^D$, followed by Cartier duality, shows that the $m$th $q$-Verschiebung kills $G$ as well. The duality used here is \citet[Proposition~5.19 and proof, pp.~79--80]{EGM}. Over a perfect field the \'etale, multiplicative and local-local factors are canonical \citep[Lemma~4.44, Proposition~4.45 and the following paragraph, pp.~67--68]{EGM}, so these statements apply to their factors in an abelian variety's torsion.

\begin{proof}[Proof of Theorem~\ref{thm:finite-field}]
First suppose that the pair $(A,M)$ is defined over $\FF_q$, where $q=2^f$, and that $a=a(M)>0$. All line-bundle roots below are taken after algebraic closure. Let $\pi$ be the $q$-Frobenius of this model and let $v$ be its Verschiebung, so that
\[
                     \pi v=v\pi=[q].
\]
These identities and their iterates are \citet[Proposition~5.20 and \S5.21, p.~80]{EGM}. Put
\[
 m=\left\lceil\frac{2a}{f}\right\rceil,
 \qquad u=\pi^m+v^m,\qquad w=\pi^m-v^m.
\]
The local-local factor $G$ of $A[2^a]$ has order $2^{2a}$, and both $G$ and $G^D$ are connected. Since $q^m\ge2^{2a}$, the preceding coordinate-algebra argument proves
\begin{equation}\label{eq:finite-kills}
                        \pi^mG=v^mG=0,
                        \qquad uG=0.
\end{equation}
This argument takes place on the finite-field model. It does not identify an arbitrary scalar extension of a relative Frobenius with the absolute power map on all scalar coefficients.

Because $\lambda_M$ is defined over $\FF_q$, Frobenius--Verschiebung duality gives
\[
 \pi^\vee\lambda_M=\lambda_Mv,
 \qquad v^\vee\lambda_M=\lambda_M\pi,
 \qquad u^\vee\lambda_M=\lambda_Mu;
\]
see \citet[Proposition~7.34 and proof, p.~111]{EGM}. On the rank-one \'etale Tate module, $\pi$ acts by a unit $z\in\mathbf Z_2^\times$ and $v$ by $q/z$. Thus $u$ acts by the unit
\[
                            z^m+q^m/z^m.
\]
Consider the actual line bundle
\begin{equation}\label{eq:finite-L}
                         L=u^*M\otimes M^{2^a}.
\end{equation}
A pullback of an ample line bundle is nef, and its tensor product with an ample bundle is ample. Hence $L$ is ample. Its polarization is
\[
 \lambda_L=u^\vee\lambda_Mu+2^a\lambda_M
           =\lambda_M(u^2+[2^a]).
\]
The ordinary Tate scalar has valuation $a$, because
$(z^m+q^m/z^m)^2+2^a$ is odd. Consequently $a(L)=a$.

The homomorphism $\lambda_L$ kills the whole group scheme $A[2^a]$. On its ordinary factor, $\lambda_M$ already kills this level by Lemma~\ref{lem:p10blocks}; on its local-local factor the first summand vanishes by \eqref{eq:finite-kills}. The second summand $2^a\lambda_M$ vanishes at this level on every factor. The two-block criterion therefore gives $b(L)\ge a$. The actual-root assertion of Theorem~\ref{thm:p10bounds} supplies an ample bundle $N$ with
\begin{equation}\label{eq:finite-root}
                         N^{2^a}\simeq L,
                         \qquad a(N)=0.
\end{equation}

We next prove the explicit estimate. Numerical equivalence of line bundles is written additively in $\NS(A)\otimes\mathbf Q$. Finite-field Frobenius gives $(\pi^m)^*[M]=q^m[M]$. Applying $(v^m)^*$ and using $\pi^mv^m=[q^m]$ gives $(v^m)^*[M]=q^m[M]$, since $[q^m]^*[M]=q^{2m}[M]$. The multiplication formula is \citet[Corollary~2.12 and proof, p.~21]{EGM}. The numerical parallelogram identity, a consequence of the theorem of the cube, now gives
\begin{equation}\label{eq:finite-parallelogram}
                  u^*[M]+w^*[M]=4q^m[M].
\end{equation}
Here $[-1]^*[M]=[M]$ numerically; the requisite pullback identity and its proof are \citet[Theorem~2.7 and Corollary~2.8, pp.~19--20]{EGM}. Combining \eqref{eq:finite-L}--\eqref{eq:finite-parallelogram} yields
\begin{equation}\label{eq:finite-nef}
 \left(1+\frac{4q^m}{2^a}\right)[M]-[N]
                  =\frac{1}{2^a}w^*[M],
\end{equation}
which is nef. For every integer $e>1+4q^m/2^a$, the class of $M^e\otimes N^{-1}$ is the sum of a positive multiple of the ample class $[M]$ and the nef class in \eqref{eq:finite-nef}. It is therefore ample. Theorem~\ref{thm:ample-difference} proves the stated defect and relative rank.

Finally suppose only that $A/k$ is the scalar extension of an abelian surface over a finite field. Choose an ample line bundle on that projective finite-field model, enlarging the finite field if necessary, and call it $P$. If $a(P)=0$, it is already the desired bundle. Otherwise the construction just given, with $P$ in place of $M$, produces an ample ordinary-index-zero root over the algebraic closure of that finite field. The bundle and its root isomorphism descend to a finite extension by finite presentation; base changing them to $k$ proves the first assertion. For an arbitrary ample $M$ on $A/k$ with $a(M)>0$, the previously constructed $N$ makes $M^e\otimes N^{-1}$ ample for all sufficiently large $e$. Theorem~\ref{thm:ample-difference} again gives the asserted eventual rank. The numerical bound uses a finite-field model of the pair $(A,M)$ and is asserted only under that additional hypothesis.
\end{proof}

The argument makes the ordinary polarization block invertible. It does not assert that $\chi(N)$ is odd: the local-local index may remain positive. Its rank conclusion is eventual, and the numerical threshold is only sufficient. Neither the unclassified smaller powers nor the uniform first-variation determinant are settled by this construction.
